\section{Goal-Conditioned 分层策略}

\subsection{高层子目标与低层执行}

高层策略输出连续/离散子目标 \(g\)，低层策略以 \((s, g)\) 为 conditioning 进行动作采样。对于 continuous case，可把子目标定义为要素的高维 embedding；对于 discrete case，可直接当作技能 id。

\begin{table}[H]
\centering
\begin{tabular}{p{3cm}p{5cm}p{4cm}}
\toprule
方法 & 核心机制 & 典型处理 \\
\midrule
HIRO & Off-policy correction & 高层每 $c$ 步重设目标 \\
HAC & Recursive hindsight relabeling & 高层 reward 直接来自子目标 \\
FuN & Predictive subgoal representations & 用 predictive model 衡量 subgoal 可达性 \\
\bottomrule
\end{tabular}
\caption{常见 goal-conditioned 分层方法对比。}
\end{table}

\begin{knowledgebox}{子目标设计考量}
\begin{itemize}
    \item 可达性：子目标必须在低层 horizon 内可实现；
    \item 可辨识性：低层策略需要知道何时完成子目标；
    \item 可复用性：设计通用子目标，避免高层频繁创造新目标；
    \item 稳定性：高层子目标分布要有限，靠 hindsight relabel 控制 drift。
\end{itemize}
\end{knowledgebox}

\subsection{训练挑战与反馈环}

层次 RL 的训练常伴随 non-stationarity：高层策略面对不断变化的低层行为，低层策略又依赖动态子目标。讲者建议使用双轨 replay buffer（分别存 high/low）并写 log 记录 high-level goal hash 与 associated reward。

\begin{warningbox}{梯度传播与回放冲突}
\begin{itemize}
    \item low-level replay 可能过拟合早期 high-level plan；
    \item high-level reward 需要依据 low-level completion 进行 hindsight relabel，否则梯度噪声过大；
    \item 使用 off-policy correction（如 HIRO）可校正 high-level update 的 bias。
\end{itemize}
\end{warningbox}

\subsection{Hindsight relabeling 与数据重用}

Hindsight relabeling 把 high-level plan 中未达成的子目标，用 low-level 成功的状态“重写”为新的 pseudo-goal，从而让 replay buffer 提供更多正样本。实施时需通过 relabel scheduler 控制频率，避免过度 rewrite 破坏 original objective；所有 relabeled sample 的 metadata 应写入 prompt log 以便 audit。

\begin{importantbox}{Hindsight relabel 的实施步骤}
\begin{enumerate}
    \item 收集一条 episode 的高层 goal + 低层 trajectory；
    \item 选择一个状态作为 relabel goal；
    \item 重新计算 reward 并更新 high-level Q；
    \item 将 relabeled sample 的 metadata（goal hash + timestamp）写入 instrumentation；
\end{enumerate}
\end{importantbox}

\subsection{本章小结}

Goal-conditioned 分层策略把高层像指挥官一样设子目标，低层视作执行者；训练难点在于频繁的环境/目标 drift，需要 hindsight+off-policy guardrail。

